\section{Blockchain Integration Architecture}
\label{sec:blockchain}

\subsection{Distributed Systems Property}

\paragraph{Why the blockchain is structurally necessary.}
Classical self-stabilization~\cite{dijkstra_self_stab,dolev_self_stab} requires a
\emph{stabilizing medium}: shared registers that processes can read and write to
coordinate recovery.
Under crash faults, standard shared memory suffices because a crashed process
simply stops writing.
Under Byzantine faults (Assumption~1), however, a faulty process can write
\emph{arbitrary} values to any shared register, corrupting the stabilizing medium
itself and breaking the monotonic progress argument.
No classical shared-memory self-stabilization protocol can tolerate this: the
stabilizing medium is precisely the component that Byzantine faults attack.

The blockchain resolves this gap structurally.
The $n$ FL clients (Assumption~1) and the $m$ blockchain validators (Assumption~4) are
\emph{distinct} populations: FL clients tolerate $f < n/2$ Byzantine faults; validators
tolerate a separate budget of $f_v < m/3$ faulty validators (PBFT~\cite{bft_consensus}).
Once round $t$ is finalized, $\mathrm{ckpt}_t = H(w_{t+1} \| t \| |\mathcal{T}_t|)$
cannot be altered or erased.
\emph{Concretely, without this substrate, Byzantine clients can replay previously
rejected gradients in later rounds (the coordinator has no immutable record of which
rounds were legitimately closed), forcing the trajectory to regress to a corrupted
state---an attack no in-memory aggregator prevents under the Byzantine fault model.}
The blockchain closes this gap: Byzantine clients cannot fabricate or rewrite the
on-chain history (Assumption~4), so each legitimately committed round advances the
trajectory permanently, enabling Lemma~\ref{lem:closure}.
Concretely: shared registers tolerate only crash faults and permit overwrites; the chain
tolerates Byzantine validators ($f_v{<}m/3$), enforces write-once semantics, and preserves
history (Table~\ref{tab:blockchain_vs_shmem}, Appendix~\ref{app:blockchain};
Fig.~\ref{fig:blockchain_stabilizer}, Appendix~\ref{app:experiments}).

\subsection{Round Structure}

Each training round proceeds as follows: (1) coordinator broadcasts $w^t$; (2) each
client computes sketch $S_i^t$ and submits its gradient hash to the smart contract via
\texttt{submitModelUpdate}; (3) coordinator waits for $|\mathcal{H}|$ submissions and
applies spectral filtering; (4) the aggregated model hash is written on-chain via
\texttt{finalizeRound}, committing round $t$ immutably.
Full smart contract interface specification is in Appendix~\ref{app:blockchain}.

\paragraph{Asynchronous operation.}
The synchronous round model (Assumption~1) is a simplification; our implementation
supports asynchronous submission where clients upload sketches as they become available.
Empirically, with maximum client delay $\tau_{\mathrm{max}} = 10$ rounds, the spectral
detection rate degrades by at most 12\% relative to synchronous aggregation (measured
over 100 training runs across three architectures), and the convergence rate
$O(\sigma\beta/\sqrt{T} + \beta^2/T)$ (where $\beta = f/n$) is maintained in practice.
A formal treatment of asynchronous self-stabilization under bounded delay remains
future work (see Section~\ref{sec:conclusion}).
